% =============================================================
%  第五单元  分一分与除法
% =============================================================
\Chapter{第五关\quad 分东西，不许打架}
\noindent{\small\kaishu\color{tipc}（课本第五单元：分一分与除法，P56--P71）}

\begin{blunder}
过生日分蛋糕，豆角爸爸一刀下去，一块大一块小。
小豆角当场指出来：“豆角爸爸偏心。”然后小豆角跑回房间拿了把尺子，
量了两块的高度给豆角爸爸看。那天豆角爸爸给小豆角道了歉，重新切。
\end{blunder}

\begin{gamble}{赌一把：分得匀不匀，数了才知道}
\begin{itemize}
  \item 抓一把豆子（别数），两人各分三堆，要求“一样多”。
  \item 分完互相检查。\textbf{差的少的人赢}；差得多的人\textbf{叠衣服}。
  \item \textbf{隐藏彩蛋}：分零食的时候定个规矩——\textbf{剩下的归分的人}。
        小豆角立刻就会算得极其认真，而且会主动想办法少剩一点。
\end{itemize}
\end{gamble}

\begin{bonus}
\begin{itemize}
  \item \textbf{“剩下的归分的人”}：分任何东西都定这条规矩
        （饼干、葡萄、糖果、甚至洗碗的活儿）。小豆角会瞬间变得极其精确。
  \item \textbf{分碗筷任务}：把分筷子这件事长期交给小豆角。
        每天一次“每人几根”，就是每天一道除法。
  \item \textbf{一包薯片多少片（买一包，数出来）}：全家几个人？每人几片？剩几片归谁？
        让小豆角自己算，算完真的按小豆角说的分。
\end{itemize}
\end{bonus}

\begin{unitinfo}
\textbf{课本 P56--P71}\quad \textbf{6 次课，每次 20 分钟}\quad
\textbf{建议第 14--16 周} \\[4pt]
\textbf{这一关到底在干嘛}：先学会“分得公平”，再把“分”写成除法算式。\\
\textbf{逻辑是这么走的}：\ 分东西要一样多 $\to$ 有两种问法（分几份／每份几个）
$\to$ 两种都能写成除法 $\to$ 除法不用算，\textbf{倒着想乘法口诀就行}
$\to$ 分不完的，就说“还剩几块”。\\
\textbf{家长必须知道的一条}：二年级上册\textbf{不教“余数”这个词}，
也不写“$\cdots\cdots$”。课本只说“每人 8 块，\textbf{还剩} 2 块”。
\textbf{别提前教“余数”}，下学期会正式学，现在教了必混。
\end{unitinfo}

\subsection{开局先聊个事：除法就是“分得公平”}

吃饭了，小豆角分筷子。四个人，每人两根，一共拿几根？
——这是乘法。

反过来：一把筷子 12 根，四个人平均分，每人几根？
——这就是除法。

分水果、分饼干、分糖果，还有更实在的：
\textbf{20 块钱，一个玩具 5 块，能买几个？}——也是除法。

所以除法不是新东西，它就是把“公平地分”这件事写下来。

\subsection{任务 1：什么叫“一样多”}

\begin{lesson}{1}{20}
\textbf{今天要拿下的}：小豆角懂“平均分”就是每份一样多，并且能动手做出来。

\begin{materials}
12 颗小饼干（或豆子）、3 个纸杯。
\end{materials}

\begin{enumerate}
  \item \textbf{先分得不公平}：12 颗饼干分 3 个杯子，故意分成 5、4、3。
        问：“这样分行吗？”（有的多有的少，不公平）
  \item \textbf{让小豆角重分}：要求每个杯子一样多。小豆角会一颗一颗轮着放。分完每杯 4 颗。
  \item \textbf{给这件事起个名字}：“每份一样多”，就叫\textbf{平均分}。
        \textbf{别一次讲太多，让小豆角先记住“一样多”三个字。}
  \item \textbf{两种分法都试一遍}（体会过程可以不同）：
        \begin{itemize}
          \item \textbf{1 个 1 个分}：每杯先放 1 颗，转一圈再放 1 颗……
          \item \textbf{2 个 2 个分}：每杯一次放 2 颗，两轮分完。
        \end{itemize}
        问：“分法不一样，结果一样吗？”（一样，每杯 4 颗）
  \item \textbf{让小豆角记下来}：把分的过程画出来或写出来（课本 P57 的做法：
        每只兔子先分 2 根，就画 2 个）。
        \textbf{记录比答案重要}，这是后面写除法算式的准备。
  \item \textbf{换总数}：15 根小棒平均分 3 杯 $\to$ 每杯 5 根；20 根分 4 杯 $\to$ 每杯 5 根。
        \textbf{每次都动手，不许心算。}
  \item \textbf{判断}：给几组数（3、3、3；4、4、5；2、2、2、2），
        哪些是平均分？为什么？
\end{enumerate}

\begin{saybox}
“这三杯一样多吗？小豆角怎么知道？”\\
“小豆角是 1 个 1 个分的，还是 2 个 2 个分的？结果一样吗？”\\
“15 根平均分 3 杯，小豆角怎么知道每杯 5 根？”（分完数的）
\end{saybox}

\begin{passbox}
\begin{itemize}
  \item 12、15、20 能独立平均分成 3 份或 4 份，每份一样多。
  \item 小豆角能说出“平均分就是每份一样多”。
  \item 小豆角能把自己分的过程画下来或写下来。
\end{itemize}
\end{passbox}
\end{lesson}

\subsection{任务 2：两种分法，两个意思}

\begin{lesson}{2}{20}
\textbf{今天要拿下的}：分清“分成几份”和“每份几个”——这是本章第一个坎。

\begin{materials}
12 根小棒（当苹果）、3 个纸杯（当盘子／袋子）。
\end{materials}

\begin{enumerate}
  \item \textbf{第一种：知道份数，求每份几个}（课本 P58）。
        12 个苹果，平均放 3 个盘子，每盘几个？
        $\to$ 动手分 $\to$ 每盘 4 个 $\to 12\div3=4$（个）。
  \item \textbf{第二种：知道每份几个，求份数}。
        同样 12 个苹果，\textbf{每袋装 3 个}，能装几袋？
        $\to$ 动手装 $\to$ 4 袋 $\to 12\div3=4$（袋）。
  \item \textbf{摆在一起看}：
        \begin{center}
        \small
        \begin{tabular}{@{}lll@{}}
        \toprule
        \textbf{问什么} & \textbf{算式} & \textbf{单位} \\
        \midrule
        平均分成 3 份，每份几个？ & $12\div3=4$ & 个 \\
        每 3 个一份，能分几份？ & $12\div3=4$ & 份（袋）\\
        \bottomrule
        \end{tabular}
        \end{center}
        \textbf{算式一模一样，单位不一样，意思就不一样。}
        这一句必须小豆角自己说出来。
  \item \textbf{换一组数}（课本 P62 分香蕉）：12 根香蕉。
        \begin{itemize}
          \item 平均分给 3 只猴子：$12\div3=4$（根）
          \item 每只猴子分 2 根：$12\div2=6$（只猴子）
        \end{itemize}
        \textbf{这两个算式不一样}——除以几，就是“分成几份”或“每份几个”。
  \item \textbf{配对练习}：四个问题，只有两个算式，让小豆角连：
        \begin{center}
        \small
        \begin{tabularx}{\linewidth}{@{}p{7.2cm}X@{}}
        \toprule
        \textbf{问题} & \textbf{算式} \\
        \midrule
        18 颗糖平均分给 6 人，每人几颗？ & $18\div6=3$ \\
        18 颗糖，每人 3 颗，能分给几人？ & $18\div3=6$ \\
        18 颗糖，每袋 6 颗，能装几袋？ & $18\div6=3$ \\
        18 颗糖平均装 3 袋，每袋几颗？ & $18\div3=6$ \\
        \bottomrule
        \end{tabularx}
        \end{center}
        小豆角会觉得奇怪：四道题只有两个算式。\textbf{正好用来说清“算式一样，意思可以不一样”。}
\end{enumerate}

\begin{saybox}
“这两个算式一模一样，意思一样吗？”\\
“小豆角怎么知道单位是‘个’还是‘份’？”（看问的是什么）\\
“每袋装 3 个，能装几袋——用除法还是用减法？”（都行，除法快）
\end{saybox}

\begin{pitbox}
\textbf{最典型的错}：两种分法混成一种。看到 12 和 3 就写 $12\div3=4$，
但问“能装几袋”时答“4 个”。\\
\textbf{解法}：\textbf{每次都要求小豆角先说“这是什么单位”}，写出来。
单位错了，说明没懂，别往下走。
\end{pitbox}

\begin{passbox}
\begin{itemize}
  \item 看情境图，小豆角能说出是哪一种分法，算式和单位都对。
  \item 小豆角能解释“算式一样、意思不一样”是怎么回事。
\end{itemize}
\end{passbox}
\end{lesson}

\subsection{任务 3：除法算式长什么样 + 用口诀求商}

\begin{lesson}{3}{20}
\textbf{今天要拿下的}：认识被除数、除数、商；不摆小棒也能用口诀求出商。

\begin{materials}
白纸、笔、附录 A 的乘法卡。
\end{materials}

\begin{enumerate}
  \item \textbf{认名字}（在白板上写个大的）：
        \begin{center}
        \begin{tikzpicture}[font=\small]
          \node (expr) at (0,0) {\Large $12\quad\div\quad3\quad=\quad4$};
          \node[below=4mm of expr, xshift=-1.25cm, text=mainc] {\scriptsize 被除数};
          \node[below=4mm of expr, xshift=0.42cm, text=mainc] {\scriptsize 除号};
          \node[below=4mm of expr, xshift=1.35cm, text=mainc] {\scriptsize 除数};
          \node[below=4mm of expr, xshift=2.75cm, text=mainc] {\scriptsize 商};
        \end{tikzpicture}
        \end{center}
        \textbf{读作}“12 除以 3”。
  \item \textbf{回到意思}：“12 是被分的总数，3 是分成几份（或每份几个），4 是结果。”
  \item \textbf{不摆小棒也能算}（课本 P64 小熊开店）：
        \begin{center}
        \small
        \begin{tabular}{@{}lll@{}}
        \toprule
        \textbf{题} & \textbf{想哪句口诀} & \textbf{答案} \\
        \midrule
        $20\div5$（20 元买 5 元的车） & 四五二十 & 4（辆）\\
        $36\div9$（36 元买 9 元的玩具） & 四九三十六 & 4（个）\\
        $24\div6$ & 四六二十四 & 4 \\
        $24\div4$ & 四六二十四 & 6 \\
        $35\div5$ & 五七三十五 & 7 \\
        \bottomrule
        \end{tabular}
        \end{center}
        \textbf{关键发现}：$24\div6$ 和 $24\div4$ \textbf{用的是同一句口诀}
        “四六二十四”。让小豆角自己说出来。
  \item \textbf{乘除一家四口}（课本 P65）：
        \begin{center}
        $6\times4=24$ \qquad $4\times6=24$ \qquad $24\div6=4$ \qquad $24\div4=6$
        \end{center}
        让小豆角观察：“乘法里的\textbf{积}，到除法里当了\textbf{被除数}。”
        \textbf{四个算式、一句口诀、一家人。}
  \item \textbf{做八道}（每题都要报口诀）：
        $18\div6$、$18\div3$、$20\div4$、$20\div5$、$27\div3$、$27\div9$、$16\div4$、$16\div8$。
  \item \textbf{讲个故事}（课本 P65 第 4 题）：挑一组算式（如 $4\times5=20$、$20\div4=5$），
        让小豆角讲一件事：“4 盒饼干，每盒 5 块，一共 20 块；20 块饼干平均放 4 盒，每盒 5 块。”
        \textbf{讲错也没关系，讲的过程就是在理解。}
\end{enumerate}

\begin{saybox}
“12 除以 3，12 是总数，3 是什么？”\\
“$24\div6$ 想哪句口诀？那 $24\div4$ 呢？”（同一句）\\
“乘法里的积，到除法里叫什么？”（被除数）
\end{saybox}

\begin{passbox}
\begin{itemize}
  \item 能说出被除数、除数、商，会说“12 除以 3”。
  \item \textbf{八道除法全用口诀直接答出来，不再摆小棒。}
  \item 能说出一句口诀对应的四个算式。
\end{itemize}
\end{passbox}
\end{lesson}

\subsection{任务 4：“是几倍”要用除法}

\begin{lesson}{4}{20}
\textbf{今天要拿下的}：分清“求几倍是多少（乘法）”和“求是几倍（除法）”。

\begin{materials}
纸、笔、彩笔。
\end{materials}

\begin{enumerate}
  \item \textbf{先复习乘法的那个方向}：蜜蜂 9 只，蝴蝶是蜜蜂的 3 倍，
        蝴蝶几只？$\to 9\times3=27$（只）。
        \textbf{知道一份和倍数，求总量 $\to$ 乘法。}
  \item \textbf{反过来}（课本 P67 花园）：蝴蝶 4 只，蜻蜓 12 只。
        蜻蜓是蝴蝶的几倍？
        $\to$ 圈一圈：12 里面有 3 个 4 $\to 12\div4=3$。\\
        \textbf{“几倍”就是“里面有几个”，所以用除法。}
  \item \textbf{第三种}：蜻蜓 12 只，是蚂蚁的 3 倍，蚂蚁几只？
        $\to$ 把 12 平均分 3 份，一份就是蚂蚁数 $\to 12\div3=4$（只）。
  \item \textbf{三种问法摆一张表，贴桌上}：
        \begin{center}
        \small
        \begin{tabularx}{\linewidth}{@{}p{6.4cm}p{3.4cm}X@{}}
        \toprule
        \textbf{问什么} & \textbf{用什么} & \textbf{例} \\
        \midrule
        求“是多少”（知道一份和倍数） & 乘法 & 9 只的 3 倍 $=9\times3=27$ \\
        求“是几倍”（知道两个量） & 除法 & 12 里有几个 4：$12\div4=3$ \\
        求“一份”（知道总数和倍数） & 除法 & 12 是蚂蚁的 3 倍：$12\div3=4$ \\
        \bottomrule
        \end{tabularx}
        \end{center}
  \item \textbf{练六道，每题先判断用乘还是除}：
        \begin{enumerate}[label=（\arabic*）,itemsep=2pt]
          \item 小轿车 6 人，大客车是小轿车的 4 倍，大客车几人？
          \item 小明 8 张卡片，小红 24 张，小红是小明的几倍？
          \item 妈妈 36 岁，是小明的 4 倍，小明几岁？
          \item 一支笔 3 元，一支钢笔 18 元，钢笔是笔的几倍？
          \item 树上有 7 只鸟，地上 21 只，地上是树上的几倍？
          \item 一箱牛奶 24 盒，平均分给 6 人，每人几盒？
        \end{enumerate}
        答案：（1）$6\times4=24$（2）$24\div8=3$（3）$36\div4=9$（4）$18\div3=6$
        （5）$21\div7=3$（6）$24\div6=4$。
\end{enumerate}

\begin{saybox}
“这题是求‘是多少’还是求‘是几倍’？”\\
“12 里面有几个 4？怎么数最快？”（4、4、4——3 个）\\
“为什么求几倍用除法？”（因为要看里面有几个）
\end{saybox}

\begin{passbox}
\begin{itemize}
  \item 三种问法能分清，六道题全对。
  \item 小豆角能说出“求几倍就是看里面有几个”。
\end{itemize}
\end{passbox}
\end{lesson}

\subsection{任务 5：分不完怎么办}

\begin{lesson}{5}{20}
\textbf{今天要拿下的}：平均分还有剩的时候，能说清“每人几块、还剩几块”。

\begin{materials}
50 根小棒（当糖）、4 个纸杯、6 个纸杯。
\end{materials}

\begin{enumerate}
  \item \textbf{先分}：50 块糖平均分给 4 个人。1 块 1 块分、2 块 2 块分、
        5 块 5 块分，都试一遍 $\to$ 每人 12 块，剩 2 块。
  \item \textbf{把过程记成表}（课本 P60 的办法）：
        \begin{center}
        \small
        \begin{tabular}{@{}lcccc@{}}
        \toprule
        & 第 1 次 & 第 2 次 & 第 3 次 & 合计 \\
        \midrule
        每人分到 & 5 & 5 & 2 & 12 \\
        分掉 & 20 & 20 & 8 & 48 \\
        还剩 & 30 & 10 & 2 & 2 \\
        \bottomrule
        \end{tabular}
        \end{center}
        \textbf{表格把分的过程留下来了}，小豆角能看见自己分了几轮。
  \item \textbf{换成 6 个人}：50 块分 6 人 $\to$ 每人 8 块，剩 2 块。
        \textbf{用乘法检查}：$6\times8=48$，比 50 少 2，所以剩 2。
        \textbf{这一步很重要：除法用乘法验算。}
  \item \textbf{说法要规范}：答案写成“每人 8 块，还剩 2 块”。
        \textbf{不写“余数”，不写“$\cdots\cdots$”。}
  \item \textbf{体会“多分几次就分完”}：50 块分 4 人，1 块 1 块分要 12 轮，
        5 块 5 块分只要 3 轮。\textbf{数大的时候每次多分几个，次数就少。}
  \item \textbf{生活里的四道}（课本 P60--P61）：
        \begin{itemize}
          \item 27 粒米，每只蚂蚁一次搬 3 粒，几只蚂蚁一次搬完？$27\div3=9$（只）
          \item 24 棵树苗，每行栽 4 棵，能栽几行？$24\div4=6$（行）
          \item 24 棵树苗平均分 3 个组，每组几棵？$24\div3=8$（棵）
        \end{itemize}
\end{enumerate}

\begin{saybox}
“还剩几块？小豆角怎么知道？”\\
“用乘法怎么检查？”（$6\times8=48$，48 加 2 才是 50）\\
“1 块 1 块分和 5 块 5 块分，结果一样吗？哪个快？”
\end{saybox}

\begin{pitbox}
\textbf{如果小豆角问“剩 2 块还能不能再分？”——非常好，说明小豆角在想。}
回小豆角：“2 块不够 6 个人一人一块了，所以不能再分，就剩着。”
\textbf{别讲“余数一定要小于除数”这种话}，二年级上册不需要，
讲了只会让小豆角多背一条用不上的规矩。
\end{pitbox}

\begin{passbox}
\begin{itemize}
  \item 50 平均分 6 人，能正确说“每人 8 块，还剩 2 块”。
  \item 会用乘法验算。
  \item 能填完“分的过程”那张表。
\end{itemize}
\end{passbox}
\end{lesson}

\subsection{任务 6：火星探险（游戏课）}

\begin{lesson}{6}{20}
\textbf{今天要拿下的}：在游戏里把乘、除、加减全用一遍。

\begin{playbox}{火星探险（课本 P71）}
\textbf{要准备}：骰子 1 颗、棋子每人 1 枚、下面这张棋盘。\\
\textbf{规则}：\\
① 轮流掷骰子，掷出几，就把棋子放到\textbf{数字是那个数}的格子（掷 3 就放 3 号格）；
那格被占了就重掷。\\
② 按那格上的\textbf{指令}算出结果，走到结果对应的数字那格。\\
③ 谁先到 50 号“找到矿石”谁赢；也可以比谁走的步数少。
\end{playbox}

\begin{center}
\small
\begin{tabularx}{\linewidth}{@{}p{0.9cm}Xp{0.9cm}Xp{0.9cm}X@{}}
\toprule
\textbf{格} & \textbf{指令} & \textbf{格} & \textbf{指令} & \textbf{格} & \textbf{指令} \\
\midrule
1 & 找到火星地图 & 20 & 遇到沙暴，退到 $62-47$ & 40 & 穿过岩石，到 5 的 6 倍 \\
2 & 到 $24\div4$ & 25 & 看到火星人，到 $26+24$ & 45 & 穿过大峡谷，到 $71-26$ \\
3 & 与指挥部联系，到 $81-41$ & 30 & 战胜风暴，到 $5\times7$ & 50 & 找到矿石（终点）\\
4 & 获得氧气，到 $45-20$ & 35 & 找到脚印，到 $17+18$ & & \\
5 & 遇到风暴，退到 $85-75$ & & & & \\
6 & 温度太高，退到 3 的 5 倍 & & & & \\
10 & 找到食物，到 $25\div5$ & & & & \\
15 & 看到化石，到 $12+8$ & & & & \\
\bottomrule
\end{tabularx}
\end{center}

\noindent\textbf{核对答案}：2 号 $\to6$；3 号 $\to40$；4 号 $\to25$；5 号退到 10；
6 号退到 15；10 号 $\to5$；15 号 $\to20$；20 号退到 15；25 号 $\to50$；
30 号 $\to35$；35 号原地；40 号 $\to30$；45 号原地。

\begin{tipbox}
\textbf{玩的时候豆角爸爸做三件事}：\\
① 每句指令都让小豆角自己算，豆角爸爸只说“对／不对”；\\
② 算错了别马上改，让小豆角走到错的那格上，后面小豆角自己会发现；\\
③ 玩两轮。第一轮重算对，第二轮开始比步数——这时小豆角会主动想“怎么走更快”。
\end{tipbox}
\end{lesson}

\subsection{翻车现场（对着看）}

\begin{center}
\small
\begin{tabularx}{\linewidth}{@{}p{3.8cm}p{4.0cm}X@{}}
\toprule
\textbf{小豆角写出来的样子} & \textbf{其实是想错了什么} & \textbf{豆角爸爸回一句什么} \\
\midrule
两种分法混成一种 & 没分清“分几份”和“每份几个” & 先问单位：是“个”还是“份／袋”？ \\
\addlinespace[2pt]
$24\div6$ 和 $24\div4$ 用两句口诀 & 不知道一句能管两道除法 & 都回到“四六二十四”。\\
\addlinespace[2pt]
求几倍用了乘法 & 三种问法没分开 & 看那张三行对照表。\\
\addlinespace[2pt]
“还剩”写成“余数” & 提前学了二年级下册的内容 & 这学期只说“还剩几块”。\\
\addlinespace[2pt]
除法不验算 & 没有检查习惯 & 用乘法验：$6\times8=48$。\\
\addlinespace[2pt]
除号写得不像除号 & 符号没练够 & 横线上下各一个点，写格子正中。\\
\bottomrule
\end{tabularx}
\end{center}

\subsection{抽查：随机开枪（10 题）}

\begin{enumerate}[label=\arabic*.,itemsep=4pt]
  \item $12\div3=$ \blank[1.4cm] \qquad 2. $18\div6=$ \blank[1.4cm] \qquad
  3. $24\div4=$ \blank[1.4cm]
  \item $35\div5=$ \blank[1.4cm] \qquad 5. $27\div9=$ \blank[1.4cm] \qquad
  6. $16\div8=$ \blank[1.4cm]
  \item $20\div5=4$ 里，被除数是 \blank[1.2cm]，除数是 \blank[1.2cm]，商是 \blank[1.2cm]。
  \item 18 个苹果平均放 3 个盘子，每盘 \blank[1.4cm] 个。算式 \blank[2.6cm]。
  \item 20 个苹果每 4 个装一袋，能装 \blank[1.4cm] 袋。算式 \blank[2.6cm]。
  \item 50 块糖平均分 6 人，每人 \blank[1.2cm] 块，还剩 \blank[1.2cm] 块。
\end{enumerate}

\begin{ansbox}
1. 4\quad 2. 3\quad 3. 6\quad 4. 7\quad 5. 3\quad 6. 2\quad
7. 20、5、4\quad 8. 6；$18\div3=6$（个）\quad 9. 5；$20\div4=5$（袋）\quad
10. 8；2 \\
\textbf{第 8、9 题单位写错}：两种分法还没分开，回任务 2。
\end{ansbox}

\subsection{课后野路子}

\begin{itemize}
  \item \textbf{分碗筷}：吃饭前让小豆角分筷子（每人 2 根）——这就是“每人几个”。
  \item \textbf{分水果}：买回来平均分给家人，每次要求小豆角说出算式和单位。
  \item \textbf{扑克牌}：抽两张，大的除以小的（除不尽就换牌），要求报出口诀。
  \item \textbf{买东西}：36 块能买几个 9 块的本子？（配合后面购物活动）
  \item \textbf{分零食}：一包饼干 24 片，三个人怎么分？分完让小豆角算每人几片、剩几片。
\end{itemize}
